We studied the strategic classification problem where a strategic sender of data
has a non-uniform preferences over labels that a receiver can assign. We analyzed two versions of this problem: a decision 
problem and a learning problem. For the decision problem, we showed that expanding the classifier’s action space by either introducing a 
penalty label that ignores certain feature vectors or introducing a randomized classification rule, 
significantly improves classification accuracy against a strategic sender. Our results show that an optimal deterministic classifier equipped with a penalty label makes fewer errors than one without a 
penalty label. Furthermore, the optimal stochastic classifier makes fewer errors than the optimal 
deterministic classifier with a penalty label.
For the learning problem, we show that incorporating the sender's best response structure can enhance the 
learning 
process. We demonstrated this by introducing and comparing two algorithms, vanilla ERM and strategy-aware ERM, which 
integrate game theory and 
learning theory in distinct ways. Vanilla ERM is preferable for simpler hypothesis classes and strategy-ware ERM is 
preferable for complex hypothesis classes. This opens avenues for exploring more nuanced combinations of game theory and learning, 
potentially leading to more powerful strategic learning algorithms. While this chapter is limited to the assumption of a 
scalar feature space or a vector feature space with a linear class boundary, and of simple utility and cost functions, 
the structural insights it provides offer guidance for understanding strategic classification in more complex settings.

